Пусть L — оператор задачи Штурма-Лиувилля. Пусть, далее, \(y_1(x)\) и \(y_2(x)\) — решения однородного уравнения \(Ly = 0\), удовлетворяющие граничным условиям на правом и левом конце соответственно. Пусть, наконец, \(\lambda = 0\) не является собственным значением задачи Штурма-Лиувилля. Тогда

\begin{enumerate}[label=\arabic*)]
    \item такие функции линейно зависимы
    \item такие функции единственны
    \item такие функции линейно независимы
    \item общее решение уравнения \(Ly = 0\) имеет вид \(y(x) = C_1 y_1(x) + C_2 y_2(x)\)
\end{enumerate}

% Сложность: 2
% Ответ: 3, 4
